\section{Fullauto: запуск, математика и проверка
результата}\label{fullauto-ux437ux430ux43fux443ux441ux43a-ux43cux430ux442ux435ux43cux430ux442ux438ux43aux430-ux438-ux43fux440ux43eux432ux435ux440ux43aux430-ux440ux435ux437ux443ux43bux44cux442ux430ux442ux430}

Руководство к объединённому коду для подготовки MBS-fast 2.0, 26 октября
2026. Источники поведения: \nolinkurl{src/FullAuto*.inc},
\nolinkurl{src/FullAuto.cpp}, \nolinkurl{src/CliOptions.cpp}. Проценты в
примерах относятся к численной погрешности внутри модели PO или GO.

\subsection{Самый короткий
запуск}\label{ux441ux430ux43cux44bux439-ux43aux43eux440ux43eux442ux43aux438ux439-ux437ux430ux43fux443ux441ux43a}

\noindent\begin{minipage}{\linewidth}
\begin{verbatim}
make -C cpu -j2
cpu/bin/mbs_po_mpi --method po --backend cpu \
  --particle 1 5 2 --ri 1.3116 0 --wavelength-um 0.532 \
  --fullauto 0.02 --output results/column_5um
\end{verbatim}
\end{minipage}

\texttt{po} --- физическая оптика; \texttt{cpu} --- вычисление на
процессоре. \texttt{-\/-particle\ 1\ L\ D} задаёт правильный
шестигранный столбик: L --- длина вдоль его оси, D --- диаметр
основания. Здесь L=5 и D=2 мкм. \texttt{-\/-ri\ REAL\ IMAG} задаёт
комплексный показатель преломления. Длина волны и координаты частицы
должны иметь одинаковые единицы. \texttt{0.02} --- допуск 2\%. Каталог
\nolinkurl{-\/-output} хранит матрицы, параметры и восстановление.

Не нужно добавлять Sobol, число φ, стартовые ориентации, коэффициенты
аналитической модели и отсечение только ради включения автоматики:
контроллер задаёт их сам. Это один процесс контроллера; не запускайте
несколько MPI-контроллеров в одну папку.

Частица из файла с максимальным габаритом Dmax=5 мкм:

\noindent\begin{minipage}{\linewidth}
\begin{verbatim}
cpu/bin/mbs_po_mpi --method po --backend cpu \
  --particle-file examples/cube.particle --resize-dmax-um 5 \
  --ri 1.3116 0 --wavelength-um 0.532 \
  --fullauto 0.02 --output results/file_5um
\end{verbatim}
\end{minipage}

Dmax --- максимальное расстояние между вершинами. Оно отличается от
длины столбика, его диаметра основания и диаметра равновеликого шара.
Данные о симметрии в файле не считаются доказательством: сначала
используется полная область ориентаций.

\subsection{GPU, конусы и подробная
команда}\label{gpu-ux43aux43eux43dux443ux441ux44b-ux438-ux43fux43eux434ux440ux43eux431ux43dux430ux44f-ux43aux43eux43cux430ux43dux434ux430}

Сборка FP64 для V100: \texttt{make\ -C\ gpu\ fp64\ GPU\_ARCH=70\ -j2}.
Для RTX 3080 Ti используйте архитектуру 86, для RTX 4090 --- 89. Сборка
\texttt{fp32} быстрее на потребительских картах, но её арифметическую
погрешность нужно проверить отдельно: статистический интервал Sobol её
не измеряет.

\noindent\begin{minipage}{\linewidth}
\begin{verbatim}
# Вперёд: 0–25 градусов, один GPU.
CUDA_VISIBLE_DEVICES=0 gpu/bin/mbs_po_gpu_double \
  --method po --backend cuda --particle 1 100 50 \
  --ri 1.3116 0 --wavelength-um 0.532 \
  --fullauto 0.02 --fullauto-theta-range 0 25 \
  --output results/column_100_forward

# Назад: 170–180 градусов, два GPU одного контроллера.
CUDA_VISIBLE_DEVICES=0,1 gpu/bin/mbs_po_gpu_double \
  --method po --backend cuda --particle 1 100 50 \
  --ri 1.3116 0 --wavelength-um 0.532 \
  --fullauto 0.02 --fullauto-theta-range 170 180 \
  --fullauto-gpus 0,1 --output results/column_100_backward
\end{verbatim}
\end{minipage}

Номера GPU отсчитываются после применения
\nolinkurl{CUDA_VISIBLE_DEVICES}. Для двух независимых размеров
запускайте два контроллера с разными каталогами и непересекающимися
списками GPU. Флаги старого multisize не совмещаются с fullauto.

Расширенный пример намеренно фиксирует часть параметров:

\noindent\begin{minipage}{\linewidth}
\begin{verbatim}
CUDA_VISIBLE_DEVICES=0 gpu/bin/mbs_po_gpu_double \
  --method po --backend cuda --particle-file examples/cube.particle --resize-dmax-um 5 \
  --ri 1.3116 0 --wavelength-um 0.532 \
  --fullauto 0.02 --fullauto-theta-range 170 180 \
  --max-reflections 8 --phi-points 4 \
  --fullauto-initial 512 --fullauto-min-correction 256 \
  --max-orientations 1048576 --threads 4 \
  --fullauto-rotation-symmetry off --fullauto-mirror off \
  --cutoff-profile off --fullauto-analytic-controls off \
  --output results/fixed_settings
\end{verbatim}
\end{minipage}

\begin{longtable}[]{@{}
  >{\raggedright\arraybackslash}p{(\linewidth - 2\tabcolsep) * \real{0.5000}}
  >{\raggedright\arraybackslash}p{(\linewidth - 2\tabcolsep) * \real{0.5000}}@{}}
\toprule\noalign{}
\begin{minipage}[b]{\linewidth}\raggedright
Флаг
\end{minipage} & \begin{minipage}[b]{\linewidth}\raggedright
Что меняет
\end{minipage} \\
\midrule\noalign{}
\endhead
\bottomrule\noalign{}
\endlastfoot
\texttt{-\/-max-reflections\ 8} & Фиксирует целевую глубину. Автопоиск
6\ldots18 выключен. Финальная проверка может обнаружить недостаточную
глубину. \\
\texttt{-\/-phi-points\ 4} & Фиксирует число наблюдательных азимутов;
поиск φ не заменит это значение. \\
\texttt{-\/-fullauto-initial\ 512} & Начальное N на один seed, а не
окончательное N. \\
\texttt{-\/-fullauto-min-correction\ 256} & Минимальное число ориентаций
поправки на один seed. \\
\texttt{-\/-max-orientations\ 1048576} & Верхняя граница N; достижение
границы не означает сходимость. \\
\texttt{-\/-threads\ 4} & Явное число потоков; автоматический подбор не
должен его переопределять. \\
\texttt{-\/-fullauto-rotation-symmetry\ off} & Не искать сокращение
области по вращениям тела. \\
\texttt{-\/-fullauto-mirror\ off} & Не использовать зеркальное
сокращение. \\
\texttt{-\/-cutoff-profile\ off} & Фиксировать расчёт без отсечения;
автоматический balanced с поправкой не включается. \\
\texttt{-\/-fullauto-analytic-controls\ off} & Вычислять полную PO без
аналитических контрольных компонент. \\
\end{longtable}

Обычно короткая команда удобнее. Расширенный вариант нужен для
воспроизводимых сравнений и проверки одной настройки. Явное ограничение
имеет приоритет, но несовместимые режимы должны завершаться ошибкой, а
не молча игнорироваться.

\subsection{Что означает требуемая
точность}\label{ux447ux442ux43e-ux43eux437ux43dux430ux447ux430ux435ux442-ux442ux440ux435ux431ux443ux435ux43cux430ux44f-ux442ux43eux447ux43dux43eux441ux442ux44c}

M --- матрица Мюллера 4×4; она переводит падающий вектор Стокса в
рассеянный. θ --- угол рассеяния: 0° вперёд, 180° назад. Для каждого
проверяемого θ контролируются все 16 элементов. Масштаб абсолютной
ошибки Mij --- M11 того же угла. Поэтому нуль M12 не приводит к делению
на M12. Дополнительно по умолчанию проверяются отношения Mij/M11; их
погрешность задаётся в абсолютных долях отношения.

Для каждого элемента и угла складываются его статистическая оценка, его
оценка интерполяции и фиксированные резервы. Максимум берётся
\textbf{после} этого сложения:

\[E=\max_{\theta,j}\{e_{\mathrm{stat},\theta j}
+e_{\mathrm{interp},\theta j}+r_{\mathrm{fixed}}\}.\]

Отдельно проверяются непосредственно вычисленная и полная выходная
матрицы. Нельзя складывать максимальные статистическую и
интерполяционную ошибки разных элементов и выдавать это за реализованный
критерий. Интерполяционный допуск равен EPS/4: для 2\% это 0,5\%, для
1\% --- 0,25\%. Резервы зависят от включённых проверок глубины,
симметрии, аналитики и исполнения. Фактические значения записаны в
файлах параметров и \nolinkurl{pointwise_budget_production.csv}.

95\%-интервалы здесь поточечные и статистические. Они не гарантируют
одновременно 95\% для всех углов, не доказывают отсутствие узкого пика
между проверками и не измеряют ошибку самой модели PO относительно
точного электромагнитного решения. Отключённая проверка обозначается как
непроверенная, а не как нулевая ошибка.

\subsection{Последовательность
работы}\label{ux43fux43eux441ux43bux435ux434ux43eux432ux430ux442ux435ux43bux44cux43dux43eux441ux442ux44c-ux440ux430ux431ux43eux442ux44b}

\begin{enumerate}
\def\labelenumi{\arabic{enumi}.}
\tightlist
\item
  Проверить аргументы, геометрию, единицы и доступные ресурсы. Построить
  запрошенную сетку вывода и стартовую оценку N. Если есть совместимый
  чекпоинт, восстановить его.
\item
  Подобрать целевую глубину внутренних взаимодействий, если она не
  задана. На CUDA перед этим проводится короткая настройка исполнения и
  φ для пилота; полный подбор выполняется после выбора глубины.
\item
  Проверить допустимые симметрии. Обучить аналитические контрольные
  коэффициенты и проверить их на других seed. Выбрать φ и параметры
  исполнения по стоимости достижения допуска, а не только по времени
  одного малого запуска.
\item
  Вычислять опорные и проверочные θ, восемь независимых реплик и уровни
  поправок. Сохранять завершённые блоки; использовать прежние блоки для
  совпадающих запросов.
\item
  Оценить статистику, остаток интерполяции и устойчивость префиксов.
  Уточнять ориентации или θ там, где это требуется. Шумный остаток
  сначала получает больше парных образцов; новый опорный θ добавляется
  при подтверждённой структуре.
\item
  Опубликовать промежуточную матрицу и оценки. После прохождения рабочих
  критериев выполнить включённые независимые и физические проверки и
  сохранить итог.
\end{enumerate}

\texttt{-\/-fullauto-final-checks\ off} отключает финальные
перепроверки. Это меняет объём проверенного результата, а не делает
пропущенные проверки успешными. По умолчанию их сохраняют. Такой режим
нельзя использовать с GO balanced, которому обязательна проверка ошибки
отсечения. При явном \nolinkurl{-\/-convergence-passes} число проходов
сохраняется.

\subsection{Стартовое N и глубина
n}\label{ux441ux442ux430ux440ux442ux43eux432ux43eux435-n-ux438-ux433ux43bux443ux431ux438ux43dux430-n}

N --- число ориентаций \textbf{на одну} Sobol-реплику и уровень. Восемь
реплик с N=512 означают 4096 ориентаций базового уровня; поправки имеют
свои N. n --- максимальная глубина внутренних взаимодействий в
трассировке.

По умолчанию coherent PO использует эмпирическую оценку по уже
завершённым расчётам, размеру, длине волны, диапазону θ, форме и
допуску. Это только начало, а не предсказанная гарантия сходимости. В
реализации:

\[\begin{aligned}
s={}&8.5128934+2.0264267\log_2(1+L_{\max}/(8\lambda))\\
&+0.6416260\,I(\theta_{\max}>90^\circ)
+4.0403518\,I(\Delta\theta>90^\circ)\\
&+0.8222766\,I(\text{controls off})
-2.7035215\,I(\text{не столбик типа 1}).
\end{aligned}\]

\[N_{\mathrm{forecast}}=2^s(0.01/\mathrm{EPS})^2,\qquad
N_0=\operatorname{pow2up}\!\left(\max(512,N_{\mathrm{forecast}}/16)\right).\]

Lmax --- наибольшее ребро в оценке геометрии; I --- индикатор условия (0
или 1); pow2up --- округление вверх до степени двойки. До округления
применяется предел, связанный с \nolinkurl{-\/-max-orientations}.
Коэффициенты модели получены по 234 записям, 73 уникальным когерентным
конфигурациям; перенос на новую форму проверяется последующей
сходимостью. Формула со знаком «не столбик» отражает эти данные, а не
общий физический закон для произвольных частиц.

После углового пилота эта оценка становится верхним предложением для
первого производственного раунда. По измеренному времени контроллер
выбирает степень двойки с ориентиром 1800 секунд работы на одну GPU; при
отсутствии измерений начинает с 512 на seed. Уже готовые полные префиксы
восьми seed и минимальные бюджеты поправок сохраняются.
\nolinkurl{initial_production_work.json} содержит предложенное и
выбранное число ориентаций, прогноз времени и признак наличия измерений.
Это ориентир стоимости, а не гарантия завершения за 30 минут. Решение
принимается без чтения значений производственной матрицы. Явный
\nolinkurl{-\/-fullauto-initial}, загруженная адаптивная конфигурация и
\texttt{-\/-fullauto-local\ off} сохраняют заданный старт. Следующие
раунды уточняются по ошибке; все критерии приёмки остаются прежними.
Ранняя первая оценка не означает, что расчёт сошёлся.

\texttt{-\/-fullauto-initial-model\ xi} возвращает старый ориентир
beta×gamma по мелкому дифракционному шагу; он часто завышает старт.
Явный \nolinkurl{-\/-fullauto-initial} заменяет обе модели. Для GO и
некогерентной PO применяется другой ограниченный старт, связанный с
оптической толщиной поглощения; GO без поглощения масштабно инвариантна.

Поиск n начинается с 6 и может поднять его до 18. Сравниваются n и n+2
на общих ориентациях: сначала уменьшается шум разности, затем решается,
нужно ли увеличивать n.~Старт пилота по умолчанию min(512,N0),
независимый от большого производственного N. Параметры:
\nolinkurl{-\/-fullauto-depth-start}, \nolinkurl{-\/-fullauto-depth-max},
\nolinkurl{-\/-fullauto-depth-pilot}, \nolinkurl{-\/-fullauto-depth-order}.
Неразрешённый пилот или достигнутый предел глубины не принимаются как
успех. Каждый следующий n начинает с исходного пилота. Увеличенное N
сохраняется только при восстановлении того же незавершённого кандидата.
После двух пилотов контроллер может перейти к следующему n: если средняя
разность уже выше бюджета или интервал больше удвоенного оставшегося
зазора. Во втором случае оценка затрат по закону N⁻¹ᐟ² предполагает
более чем четырёхкратное увеличение выборки. Это эвристика порядка
поиска, а не гарантия скорости сходимости Sobol. Она может выбрать
большее n.~На максимальном n это сокращение поиска выключено. Принятие
всегда требует прежней верхней границы ошибки; решение и причина
записываются как \nolinkurl{depth_pilot_decision} в
\nolinkurl{fullauto_history.jsonl}. \texttt{-\/-max-reflections\ N}
отключает поиск целевого n, но не обещает прохождение включённого
финального сравнения n/n+2.

\subsection{Усреднение Sobol и три угла
Эйлера}\label{ux443ux441ux440ux435ux434ux43dux435ux43dux438ux435-sobol-ux438-ux442ux440ux438-ux443ux433ux43bux430-ux44dux439ux43bux435ux440ux430}

Поворот задаётся R=Rz(α)Ry(β)Rz(γ). Для равномерного случайного
положения частицы нужна мера Хаара --- равномерная мера на группе
трёхмерных поворотов SO(3):

\[\langle M\rangle=\frac{1}{8\pi^2}
\int_0^{2\pi}\!d\alpha\int_0^\pi\!\sin\beta\,d\beta
\int_0^{2\pi}\!M(R,\theta)\,d\gamma.\]

Для полной области точки u,v,w из единичного куба превращаются в
cosβ=1−2u, γ=2πv, α=2πw. Равномерно раскладывать β без веса sinβ было бы
неверно. Sobol строит последовательность с равномерным покрытием куба;
вложенное перемешивание Owen с seed делает независимые рандомизированные
реплики. Это не независимые случайные точки внутри каждой реплики и не
декартова сетка Nα×Nβ×Nγ. Поэтому сравнение одного N Sobol с
произведением трёх сеточных размеров не означает такое же увеличение
точности или ускорение.

Восемь seed используются для оценки неопределённости. Пусть Yr --- целая
оценка матрицы из реплики r, включая связанные поправки. Тогда

\[\bar Y=\frac18\sum_{r=1}^8Y_r,\quad
s_Y^2=\frac17\sum_{r=1}^8(Y_r-\bar Y)^2,\quad
h_{95}=t_{0.975,7}\,s_Y/\sqrt8.\]

t --- квантиль распределения Стьюдента, около 2,365. Для Mij/M11
учитывается ковариация числителя и знаменателя и проверка разрешимости
отношения; малый шумный знаменатель не разрешается заменой на удобный
нулевой интервал. Данные половинных префиксов дополнительно проверяют
устойчивость.

N растёт степенями двойки. Уже рассчитанный префикс используется
повторно; \nolinkurl{-\/-sobol-offset} в рабочих заданиях задаёт его
продолжение. Производственные, обучающие и проверочные seed разделены.
Один seed не заменяет восемь реплик.

Sobol --- наиболее практичный текущий вариант для наших неоднородных
PO-интегралов и локального уточнения: можно наращивать префикс и
совместно считать разности уровней. Это не доказательство превосходства
над Лебедевым или разложением Вигнера на любой форме. Их преимущества
зависят от гладкости и симметрии конкретной задачи.

\subsection{Зачем φ, если уже есть
α}\label{ux437ux430ux447ux435ux43c-ux3c6-ux435ux441ux43bux438-ux443ux436ux435-ux435ux441ux442ux44c-ux3b1}

α вращает частицу вокруг направления падающего света. φ задаёт азимут
наблюдателя. При полном усреднении по α дополнительные φ могут снижать
дисперсию, сохраняя математическое ожидание. В каждой ориентации это
разные операции; для поляризационных элементов также нужен корректный
поворот базиса Стокса.

\nolinkurl{-\/-haar-alpha} добавляет третью координату Sobol. Fullauto PO
использует её сам. По мере увеличения числа ориентаций не требуется
автоматически делать φ столь же мелким, как в старой двухугловой сетке.
Но φ=1 не гарантирует минимальное время: дополнительные направления
наблюдения используют уже трассированные пучки. Контроллер выбирает
компромисс «стоимость × дисперсия» и может уточнять φ локально.
\texttt{-\/-phi-points\ N} закрепляет правило.
\texttt{-\/-sobol-seed\ 32768\ SEED} относится к ручному запуску; в
fullauto используйте \texttt{-\/-fullauto-initial\ 32768} для начального
N.

Вращательная и зеркальная симметрии относятся к телу и проверяются
отдельно. Шестикратная симметрия столбика не даёт права делить
лабораторное α или φ на шесть. Для произвольной частицы без
подтверждённой симметрии остаётся полная область.

Отказ аналитической контрольной модели не отменяет зеркальную симметрию
формы. После такого отказа fullauto заново проверяет
\texttt{mirror-gamma} на обычной PO, при выбранных φ и глубине
отражений. Если парная проверка всех 16 элементов проходит прежний
допуск \nolinkurl{epsilon/256}, зеркальное ускорение сохраняется. При
неуспехе автоматический режим использует полный диапазон γ. Явное
\texttt{-\/-fullauto-mirror\ off} сохраняется; явное
\nolinkurl{-\/-mirror-gamma} при провале проверки приводит к ошибке, а не к
скрытой смене режима. Решение записывается в
\nolinkurl{analytic_fallback_mirror.json}, численная проверка --- в
\nolinkurl{mirror_verification.json}. Проверка привязана к углам, φ,
глубине, бинарнику, режиму и коэффициентам аналитического контроля.
Контрольная точка сохраняет выбранный режим: обновление программы само
по себе не переключает уже сохранённый расчёт с полного γ на зеркальное
восстановление.

\subsection{Что вычисляется
аналитически}\label{ux447ux442ux43e-ux432ux44bux447ux438ux441ux43bux44fux435ux442ux441ux44f-ux430ux43dux430ux43bux438ux442ux438ux447ux435ux441ux43aux438}

\subsubsection{Дифракция отдельного
пучка}\label{ux434ux438ux444ux440ux430ux43aux446ux438ux44f-ux43eux442ux434ux435ux43bux44cux43dux43eux433ux43e-ux43fux443ux447ux43aux430}

Плоский многоугольный выходной пучок даёт интеграл вида

\[F(\mathbf q)=\int_A\exp(i\mathbf q\cdot\mathbf x)\,dA.\]

Для постоянной амплитуды на плоской апертуре его можно свести к сумме по
рёбрам. При малой фазе и почти нулевых знаменателях используются
устойчивые пределы и моменты. Это аналитическое вычисление дифракции
\textbf{для данной ориентации}; оно не выполняет всего ориентационного
усреднения. Комплексные амплитуды пучков складываются когерентно до
перехода к матрице Мюллера. В некогерентном режиме складываются
отдельные интенсивности: это другая физическая постановка.

\subsubsection{Аналитические контрольные
компоненты}\label{ux430ux43dux430ux43bux438ux442ux438ux447ux435ux441ux43aux438ux435-ux43aux43eux43dux442ux440ux43eux43bux44cux43dux44bux435-ux43aux43eux43cux43fux43eux43dux435ux43dux442ux44b}

Пусть C(R,θ) --- дешёвая физическая модель с известным средним μC(θ), а
M --- полная PO. Контрольная переменная даёт оценку

\[\widehat M=\overline{M-bC}+b\mu_C.\]

b --- коэффициент уменьшения дисперсии, обученный отдельно от
производственной выборки. При корректном μC вычитание и возврат среднего
сохраняют ожидание. Внутренние пути, затенение и интерференция,
отсутствующие в C, остаются в M−bC. Поэтому усреднить коэффициент
Френеля и заменить им все отражения недостаточно.

Контроли включают самовклады отражающих граней и модель тени. Для M11
используются их физические средние; поляризационные контроли и
поворотные компоненты проходят собственную проверку. Неподходящие
компоненты отключаются после проверки по полной PO: это допустимый отказ
от ускорения, а не отказ от расчёта рассеяния.

\subsubsection{Моменты Френеля и фазовые
средние}\label{ux43cux43eux43cux435ux43dux442ux44b-ux444ux440ux435ux43dux435ux43bux44f-ux438-ux444ux430ux437ux43eux432ux44bux435-ux441ux440ux435ux434ux43dux438ux435}

μ --- косинус угла падения на грань, rs(μ), rp(μ) --- коэффициенты
Френеля двух поляризаций. Отражательная способность
R(μ)=(\textbar rs\textbar²+\textbar rp\textbar²)/2 зависит от μ. Момент
--- интеграл этой функции с весом; например,

\[c_\ell=\frac{2\ell+1}{2}\int_0^1 R(\mu)P_\ell(\mu)\,d\mu,\]

где \(P_\ell\) --- полином Лежандра. В физических ядрах присутствуют
также геометрические и проекционные веса. Теорема сложения и разложение
плоской волны позволяют аналитически усреднить угловые фазовые
множители, используя сферические функции Бесселя. Общие моменты повторно
используются для многих θ. Гладкие одномерные интегралы Френеля, таблицы
радиальных ядер и часть интегралов по рёбрам остаются численными. Это
аналитико-численный контроль, не замкнутое решение всех путей.

\begin{longtable}[]{@{}
  >{\raggedright\arraybackslash}p{(\linewidth - 2\tabcolsep) * \real{0.5000}}
  >{\raggedright\arraybackslash}p{(\linewidth - 2\tabcolsep) * \real{0.5000}}@{}}
\toprule\noalign{}
\begin{minipage}[b]{\linewidth}\raggedright
Настройка
\end{minipage} & \begin{minipage}[b]{\linewidth}\raggedright
Смысл
\end{minipage} \\
\midrule\noalign{}
\endhead
\bottomrule\noalign{}
\endlastfoot
\nolinkurl{-\/-analytic-facet-average} & Явно запросить аналитический
контроль граней; fullauto умеет включать и проверять его
автоматически. \\
\texttt{-\/-analytic-shadow-control\ circular} & Контроль тени круглой
апертурой с геометрически определённой площадью. Отличие от реальной
формы остаётся в остатке. \\
\texttt{-\/-analytic-mean-cache\ FILE} & Файл кэша средних для
совпадающих параметров. В fullauto этот параметр можно оставить
автоматике. \\
\texttt{-\/-analytic-mean-threads\ N} & Потоки подготовки аналитического
среднего; явное значение фиксируется. \\
\texttt{-\/-fullauto-analytic-controls\ off} & Полная численная PO, без
этих контролей. \\
\end{longtable}

Для большой частицы подготовка среднего может стоить заметного времени и
RAM. Кэш экономит повторную подготовку; это не дополнительный физический
эффект. Режим GO и некогерентный режим не используют те же когерентные
PO-контроли.

\subsection{Поправки по глубине и
отсечению}\label{ux43fux43eux43fux440ux430ux432ux43aux438-ux43fux43e-ux433ux43bux443ux431ux438ux43dux435-ux438-ux43eux442ux441ux435ux447ux435ux43dux438ux44e}

Несколько уровней используют тождество

\[\mathbb E[M_{n_L}]=\mathbb E[M_{n_0}]
+\sum_{\ell=1}^{L}\mathbb E[M_{n_\ell}-M_{n_{\ell-1}}].\]

Это многоуровневая оценка: дорогие разности могут иметь меньшую
дисперсию, чем полная матрица, и потому требовать меньше ориентаций. Оба
слагаемых разности используют одинаковые ориентации. Где поддерживается,
общая трассировка сохраняет результаты для двух глубин. Разность берётся
между матрицами полных когерентных решений соответствующей глубины, а не
между интенсивностями отдельно отброшенных лучей. Ковариации учитываются
через полную сумму внутри каждой из восьми реплик.

По умолчанию когерентная PO использует balanced с парной поправкой:

\[\widehat M=\overline{B}_N+\overline{(O-B)}_m.\]

B --- расчёт с отсечением, O --- без отсечения. Ориентации O и B в
поправке совпадают. Физическая цель остаётся O. Небольшой m экономит
работу, пока дисперсия поправки достаточно мала; в коде используются
минимальные блоки и доля префикса. Отдельный подбор пресетов по
умолчанию больше не нужен. Время всё равно зависит от числа пучков
поправки и её дисперсии: гарантированного ускорения нет.

Явный \texttt{-\/-cutoff-profile\ balanced} фиксирует пресет и не
эквивалентен автоматической схеме с O−B. Явные пороги также имеют
приоритет. Для GO balanced применяется отдельная проверка смещения
относительно off, описанная в \nolinkurl{GO_BALANCED.md}.

\subsection{θ, интерполяция и короткие
раунды}\label{ux3b8-ux438ux43dux442ux435ux440ux43fux43eux43bux44fux446ux438ux44f-ux438-ux43aux43eux440ux43eux442ux43aux438ux435-ux440ux430ux443ux43dux434ux44b}

Выходная сетка, опорные θ и проверочные θ --- разные множества. Без
внешней сетки fullauto строит достаточно плотную сетку вывода из
размера, λ и диапазона. Прямой расчёт производится на адаптивном
подмножестве. \texttt{-\/-theta-grid-file\ FILE} фиксирует именно углы
вывода, но сам по себе не означает прямое вычисление каждой строки.

Локальная кубическая интерполяция I линейна по значениям матрицы при
фиксированных узлах. Проверяется остаток
\(\Delta=M(\theta)-I[M_{\mathrm{support}}](\theta)\) на одинаковых
ориентациях. Это позволяет сокращать общий ориентационный шум. Бюджет
остатка теперь отдельный: он не наследует максимальное N соседнего
дорогого угла. При различных рабочих φ для остатка обе стороны
используют единое правило φ. Основной расчёт сохраняет свои правила.
Выборка остатка фиксируется до чтения значений, а затем уточняется по
неуспешной проверке.

Новый опорный угол начинается с обычного стартового бюджета. Полные уже
готовые префиксы восьми seed сохраняются. Планировщик ограничивает объём
следующего раунда по измеренной стоимости; ориентир 1800 рабочих секунд
на GPU не является жёстким сроком. Неделимый префикс, трассировка или
чтение кэша могут занимать больше.

Важная проверка на реальных данных: для столбика 794,3×196,2 мкм,
λ=0,532, m=1,3116, участка 170,042--170,072°, восьми seed и N=512 на
уровень остаточная оценка составила 322,691\% при смешанных φ=4/64 и
0,0329171\% при общем φ=64. Это снижение шума \textbf{проверочного
остатка на этом участке}, не ошибка всей матрицы и не измеренное
ускорение всего расчёта. Уровни 14/16/18 считались без отсечения.

В полном DAT проверочный угол также может содержать интерполяцию, хотя
для него есть отдельный прямой расчёт. Для различения используйте
\nolinkurl{value_source} и \nolinkurl{point_type} в аннотированном CSV и
\nolinkurl{mueller_evaluated.dat}. Не заменяйте одни значения другими
постфактум без новой проверки принятого выходного оценивателя.

\subsection{Логи, результат и
восстановление}\label{ux43bux43eux433ux438-ux440ux435ux437ux443ux43bux44cux442ux430ux442-ux438-ux432ux43eux441ux441ux442ux430ux43dux43eux432ux43bux435ux43dux438ux435}

При первом восстановлении старой папки актуальный контроллер создаёт
общий индекс \nolinkurl{completed_jobs.journal}. Первый переход требует
прочитать старые метаданные; при следующих восстановлениях читается
журнал и проверяются только не вошедшие в него папки. Существующие
матрицы и маркеры завершения не переписываются. Индекс ускоряет поиск:
хеши результатов и зависимости источников всё равно проверяются до
включения сохранённого блока в оценку.

При восстановимой ошибке бюджета RAM постоянный рабочий процесс
освобождается, а незавершённый блок повторяется один раз в новом
процессе на тех же ориентациях. Частичный результат не принимается;
исходный журнал сохраняется как \texttt{stdout.memory-retry.log}. Если
сработал системный резерв RAM, контроллер ждёт его восстановления не
более 60 секунд после освобождения рабочего процесса. Ожидание и повтор
отражаются в \nolinkurl{resource_wait.json} и журнале событий. Резерв не
уменьшается. Повторная нехватка памяти или другая
вычислительная/CUDA-ошибка останавливает задачу с сохранением
завершённых чекпоинтов.

\begin{longtable}[]{@{}
  >{\raggedright\arraybackslash}p{(\linewidth - 2\tabcolsep) * \real{0.5000}}
  >{\raggedright\arraybackslash}p{(\linewidth - 2\tabcolsep) * \real{0.5000}}@{}}
\toprule\noalign{}
\begin{minipage}[b]{\linewidth}\raggedright
Файл
\end{minipage} & \begin{minipage}[b]{\linewidth}\raggedright
Назначение
\end{minipage} \\
\midrule\noalign{}
\endhead
\bottomrule\noalign{}
\endlastfoot
\nolinkurl{fullauto_status.json} & Статус, параметры и последняя
опубликованная оценка. Только \texttt{converged} означает прохождение
включённых критериев. \\
\nolinkurl{intermediate_current.json} & Ссылка на последний снимок
матрицы; подходит для просмотра незавершённого расчёта. \\
\nolinkurl{round_snapshots/} & Матрицы и точности предыдущих
опубликованных раундов. \\
\nolinkurl{parameters_accuracy.csv} & Производственные N по уровням, φ,
оценки элементов и отношений в снимке. \\
\nolinkurl{interpolation_orientation_budget.csv} & Отдельные N
проверочного остатка в снимке. \\
\nolinkurl{mueller_fullauto.dat} & Принятая или опубликованная полная
сетка; интерпретировать вместе со статусом. \\
\nolinkurl{mueller_evaluated.dat} & Прямые вычисления текущего
оценивателя. \\
\nolinkurl{mueller_fullauto_annotated.csv} & Происхождение строк и оценки
ошибок. Пустое поле не означает нулевую ошибку. \\
\nolinkurl{execution_progress.json} & Текущие PID, GPU, пакеты, времена и
число сохранённых заданий. Это не новая оценка точности. \\
\nolinkurl{fullauto_history.jsonl} & История подбора и исполнения. \\
\nolinkurl{controller_command.json} & Воспроизводимая команда. \\
\nolinkurl{jobs/}, чекпоинты и индексы & Повторное использование
завершённой работы. Сохраняйте каталог целиком. \\
\end{longtable}

\noindent\begin{minipage}{\linewidth}
\begin{verbatim}
python3 -m json.tool results/column_5um/fullauto_status.json
python3 -m json.tool results/column_5um/intermediate_current.json
python3 -m json.tool results/column_5um/execution_progress.json
\end{verbatim}
\end{minipage}

После отключения повторите исходную команду тем же бинарником и в ту же
папку. Завершённые блоки загружаются; незавершённый блок может
повториться. Чекпоинт версии 6 хранит производственный и остаточный
бюджеты, а также φ остатка. Поддержано чтение старых схем, но
совместимость бинарников проверяется отдельно. Нельзя вручную подменять
хэш, редактировать готовые матрицы или переносить только один TSV
чекпоинта без зависимых блоков.

Для нового бинарника
\texttt{-\/-fullauto-resume-compatible\ OLD\_BINARY} допустим только при
подтверждённой совместимости конкретной пары бинарников и вычислительных
профилей. Сам факт успешной сборки такой совместимостью не является. Для
нового допуска или окна используйте документированный режим
\texttt{-\/-fullauto-refine-window\ A\ B\ EPS} и сохранённую исходную
конфигурацию; не редактируйте физические параметры существующего
каталога.

Служба Linux может автоматически возобновлять процесс после
перезагрузки:

\noindent\begin{minipage}{\linewidth}
\begin{verbatim}
scripts/run_fullauto_service.sh mbs-column cpu/bin/mbs_po_mpi \
  results/column_service.log --method po --backend cpu \
  --particle 1 5 2 --ri 1.3116 0 --wavelength-um 0.532 \
  --fullauto 0.02 --output results/column_service
\end{verbatim}
\end{minipage}

Этот скрипт устанавливает пользовательскую службу systemd и запускает
её. Он управляет процессом; численные решения остаются внутри программы.

\subsection{Память, загрузка GPU и границы
метода}\label{ux43fux430ux43cux44fux442ux44c-ux437ux430ux433ux440ux443ux437ux43aux430-gpu-ux438-ux433ux440ux430ux43dux438ux446ux44b-ux43cux435ux442ux43eux434ux430}

CPU трассирует и подготавливает пучки; CUDA вычисляет дифракцию и
поддержанные части трассировки. Постоянные рабочие процессы сохраняют
CUDA-контекст и ограниченные кэши. Готовые пакеты распределяются между
GPU, проверка и сохранение данных перекрываются с дальнейшей работой.
Видеопамять занята пучками, угловыми буферами, матрицами и рабочими
областями; RAM --- геометрией, подготовкой и кэшем блоков.

\texttt{-\/-fullauto-trace-backend\ auto} сравнивает CPU, резидентную
GPU-трассировку (пучки остаются в видеопамяти) и совместную
CPU+GPU-трассировку. Принимается вариант, который быстрее минимум на 5\%
и проходит проверку всех 16 элементов и отношений к M11 на одинаковых
ориентациях. Порог расхождения --- \texttt{min(1e-7,\ EPS/1024)} для
double и \texttt{min(1e-5,\ EPS/1024)} для float. Результаты выбора
сохраняются в \nolinkurl{trace_backend_calibration.json}. При начальной
выборке меньше 1024 этот дополнительный поиск пропускается. Явный выбор
\texttt{cpu}, \texttt{resident} или \texttt{hybrid} отключает поиск.

CPU-трассировка распределяет независимые ориентации между потоками;
дерево пучков одной ориентации обрабатывается последовательно. Число
потоков подготовки нужно согласовывать с физическими ядрами, доступными
процессу через CPU affinity. Например, при выделении восьми ядер новый
расчёт запускают с \texttt{-\/-threads\ 8}. Для существующего
CUDA-расчёта без явных флагов потоков используют совместимое
восстановление и \texttt{-\/-fullauto-dispatch-threads\ 8}: меняются
пулы подготовки и аналитического среднего, исходные идентификаторы и
настройка накопления сохраняются. Этот флаг нельзя добавлять вместе с
исходными \nolinkurl{-\/-threads}, \nolinkurl{-\/-po-prepare-threads} или
\nolinkurl{-\/-analytic-mean-threads}. Фактическое число показывает строка
\texttt{PO\ preparation\ threads} в журнале работника.

Резидентный путь поддерживает выпуклые непоглощающие частицы до 32
граней, 24 вершин пучка и 18 внутренних столкновений, включая 20-гранный
дроксталл. Геометрия и оптическая фаза при подготовке остаются в double
даже в сборке float. Для частиц с более чем 16 гранями один пакет
ограничен 64 ориентациями; полное число ориентаций не сокращается.
Очередь пучков зависит от числа граней и глубины, а выделение памяти
проверяет свободную VRAM. При переполнении весь пакет пересчитывается на
CPU. Некогерентный режим поддержан: матрица Мюллера вычисляется для
каждого пучка до суммирования, без межпучковой интерференции. Поглощение
и другие неподдержанные варианты сохраняют CPU-подготовку.

Постоянный работник также запоминает консервативную оценку RAM на
ориентацию для одной физической задачи. После 128 измеренных ориентаций
он может начать следующий запрос с уже проверенного размера пакета
вместо повторного разгона от 16. Старт ограничен 256, предыдущим
измеренным размером, явным \nolinkurl{-\/-orientation-chunk} и текущей
доступной RAM с дополнительным запасом 2×. Это не кэш матриц и не
изменение выборки. Сообщение \texttt{Prepared\ memory\ reuse} показывает
оценку и стартовый пакет; кэш пучков можно независимо отключить через
\texttt{-\/-trace-cache-mib\ 0}.

В очереди совместимые угловые пакеты одного seed, префикса, n и φ идут
подряд, чтобы подготовленные пучки реже вытеснялись из ограниченного
кэша другими seeds. Выигрыш возможен, когда пучки одного задания
помещаются в кэш; объём памяти и число рассчитываемых ориентаций не
увеличиваются. Короткий пакет может оставаться после вычитания уже
сохранённых углов или разделения по n и φ.

\nolinkurl{execution_progress.json} содержит счётчики сохранённых блоков
текущей фазы и последние 16 завершённых блоков с seed, offset, N, n, φ и
границами θ. Между границами могут быть пропуски; точные углы находятся
в файлах задания. Это выполненная работа, а не новая оценка ошибки и не
процент готовности всего адаптивного раунда. Оценка точности берётся из
\nolinkurl{intermediate_current.json}.

Загрузка GPU может падать при трассировке, сборке префиксов, проверке и
I/O. 100\% загрузки не является критерием ни правильности, ни
минимального времени. Большой пакет снижает накладные расходы, но
повышает память и задержку публикации. Контроллер ограничивает память и
объём работы; OOM или исчерпание диска не являются успешным завершением.
Текущее время задания видно отдельно от возраста оценки.

\subsection{GO и extinction-only}\label{go-ux438-extinction-only}

\noindent\begin{minipage}{\linewidth}
\begin{verbatim}
cpu/bin/mbs_po_mpi --method go --backend cpu \
  --particle 1 100 50 --ri 1.3116 0 --wavelength-um 0.532 \
  --fullauto 0.02 --output results/column_go
\end{verbatim}
\end{minipage}

GO --- геометрическая оптика: α усредняется аналитически; конечные
угловые ячейки считаются напрямую, кубической PO-интерполяции нет.
Полюсные дельта-вклады сохраняются отдельно в
\nolinkurl{mueller_forward.dat} и \nolinkurl{mueller_backward.dat}; это
интегральные массы, не значения гладкой диаграммы в градусной точке.

GO mirror по умолчанию выключен. Явное включение поддерживает
непоглощающие встроенные типы 1 и 10 с допустимыми симметриями;
зеркальная пара восстанавливается через (M+PMP)/2, P=diag(1,1,−1,−1),
включая полюсные вклады. Это не простое удвоение числа симметрий.
Файловые формы и поглощение не принимаются как поддержанные.
Подробности: \nolinkurl{GO_MIRROR.md} и \nolinkurl{GO_BALANCED.md}.

\nolinkurl{-\/-extinction-only} --- отдельный CPU-режим PO с Hammersley для
выпуклой геометрии. Он вычисляет экстинкцию по оптической теореме без
угловой сетки матриц; не является fullauto и не даёт полного интеграла
рассеяния или точного альбедо. Примеры: \nolinkurl{EXTINCTION_ONLY.md}.
